\subsection{Deployment and Operations}\label{sec:deployment}
\subsubsection{Local deployment} The whole system runs locally in Docker, as three Compose projects that share one network (\code{movies-net}): \code{infra/compose.yaml} (databases, Kafka, Ollama, Keycloak), \code{backend/compose.yaml} (the four Quarkus services) and \code{frontend/compose.yaml} (the website). The infrastructure project runs three PostgreSQL~16 containers (the movies one with the pgvector extension), a single-broker Kafka in KRaft mode, Ollama with a one-shot container that downloads the two models, and Keycloak~26 (\code{start-dev --import-realm}), which imports the \code{movies} realm from \code{movies-realm.json} and keeps its state in its own database in the movies server. The movies container mounts \code{infra/init/}, whose script creates the \code{keycloak} role and database on first start. Every container has a health check (\code{pg\_isready}, Kafka's broker API, Ollama's \code{list}, the \code{/health/ready} endpoints). Because \code{depends\_on} cannot cross projects, a \code{justfile} starts them in order and waits for each to report healthy (\code{--wait}): \code{just local-first-up} also imports the \code{pg\_dump} catalogue before the backend starts, and \code{just local-up} skips the import. In containers the services run in Quarkus' production profile and read their database URLs, OIDC issuer and CORS origins from environment variables; Keycloak keeps the browser-facing address as the token issuer, while the services reach it by container name through its backchannel hostname setting.

\subsubsection{Containerization} A single multi-stage \code{backend/Dockerfile} builds any of the four Quarkus services from source (Maven in the build stage, a UBI9 OpenJDK~21 runtime in the final one), and \code{frontend/Dockerfile} builds the Next.js standalone server. The same Dockerfiles serve Compose and the cloud deployment, where \code{build-images.sh} builds the five images for \code{linux/amd64} and pushes them to Docker Hub.

\subsubsection{Kubernetes on GCP} The production target is a self-hosted cluster (managed offerings such as GKE are excluded by the specification), created entirely through IaC and driven from a management container that holds Terraform, Ansible, \code{kubectl} and \code{gcloud}. Terraform creates, in a separate stack that outlives the VMs, one data disk per worker and a reserved static IP, then a VPC and subnet, firewall rules and three Compute Engine VMs in \code{europe-west4-b} running Ubuntu 24.04: one control-plane node (\code{e2-standard-2}) and two workers (\code{e2-standard-4}); it also writes Ansible's inventory. Four Ansible playbooks follow: they prepare the nodes (containerd, kubeadm, kubelet), run \code{kubeadm init} with Calico~3.32 as pod network and join the workers; install the platform (local volumes on the workers' data disks, the Strimzi Kafka operator and kube-prometheus-stack through Helm); and deploy the application from Kubernetes manifest templates. The application namespace holds one PostgreSQL StatefulSet per service (each with a \code{postgres\_exporter} sidecar), a three-node Kafka cluster in KRaft mode (topics replicated three times), Ollama with a volume for the models and a Job that downloads them, Keycloak in production mode (realm imported from a ConfigMap), a Deployment per service with startup, readiness and liveness probes on \code{/q/health}, and NetworkPolicies that let each database accept connections only from its own service and Ollama only from the worker and the Movies service. NGINX runs on the worker that owns the static IP and is the only public entry point: it terminates TLS (a self-signed certificate unless one is provided) and routes \code{/}, \code{/api/*}, \code{/auth/} and \code{/grafana/}. Secrets are generated on the first run and kept out of git. A full deployment from nothing takes about 25 minutes.

\subsubsection{Monitoring} Prometheus and Grafana come from kube-prometheus-stack. Rather than rely on the metrics that libraries and exporters give for free, each service records its own \emph{application metrics} with Micrometer, published at \code{/q/metrics} and scraped every 15\,s: searches by whether they found anything, candidate queries by ranking (meaning, words or taste), embedding calls to Ollama and how often a request's vector was already known, the movies still to embed; profile changes by action and what the profiles database holds; recommendation requests by kind and outcome (including those answered without the model, by reusing a set made from identical inputs), the time from a request to its answer, the requests waiting and the users with the app open; and in the worker, the time each job waited in Kafka and took, plus each model call's duration and tokens, since Ollama has no metrics endpoint. Timers publish histograms up to an hour, because a CPU-bound model call far exceeds Micrometer's default 30\,s range, and totals that need a database query are counted on a schedule rather than at each scrape. A provisioned Grafana dashboard, \emph{Movies}, is built on these metrics; a last row keeps the standard ones (HTTP per API, NGINX, Kafka Exporter's consumer lag, the databases' exporters, CPU and memory per pod and node) for context. Grafana is served at \code{/grafana} behind the same address; Prometheus is not public and is reached with a port-forward. NGINX must repeat its proxy headers in the Grafana route, since it drops the inherited ones when a route sets its own; otherwise Grafana's origin check rejects every query. \TODO{figure of the dashboard under load.}

\subsubsection{Reproducibility} Step-by-step instructions (prerequisites, configuration defaults, data loading, troubleshooting) are maintained in the repository \code{README.md} and \code{infra/README.md}. Locally, \code{just local-first-up} starts the whole stack (\code{just local-up} on later runs). In the cloud, \code{just images} pushes the images, and \code{just all}, run inside the management container, runs Terraform and the three playbooks in order; \code{just destroy} removes the VMs and network but keeps the data disks and static IP, and \code{just destroy-data} removes those.
